\documentclass[a4paper,12pt]{article}
\usepackage[margin=2.5cm]{geometry}
\usepackage{hyperref}
\usepackage{microtype}

\begin{document}

\noindent
Widad Kchikach\\
Alfred Wegener Institute Helmholtz Centre for Polar and Marine Research\\
Bremerhaven, Germany\\
\texttt{widad.kchikach@awi.de}

\bigskip
\noindent
\today

\bigskip
\noindent
\textbf{Editor-in-Chief}\\
\textit{SoftwareX}

\bigskip
\noindent
Dear Editor,

\bigskip
Please find enclosed our manuscript,
\textbf{``\texttt{weatherisk}: A Python Package for Spatial Clustering
of Climate Extremes via Max-Stable Processes''}, which we are
submitting for consideration in \textit{SoftwareX}.

\bigskip
The work grew out of a practical need we kept running into during our
own research: applying spatial-extremes methods to gridded climate
output requires chaining together several technically demanding
steps---trend removal, block-maxima extraction, GEV fitting,
Fr\'{e}chet transformation, local max-stable parameter estimation,
and finally clustering---and there was no single Python tool that
handled all of them. The R implementation of Contzen, Dickhaus, and Lohmann
(``Regionalization approaches for the spatial analysis of extremal
dependence'', \textit{Extremes}, 2025,
\url{https://doi.org/10.1007/s10687-025-00519-2}) exists
and is well-validated, but porting results back and forth between R
and Python for a large HPC run is cumbersome. We wrote
\texttt{weatherisk} to close that gap.

\bigskip
The package (MIT licence, v0.1.1; archived at
\url{https://doi.org/10.5281/zenodo.20108563}) exposes the full
pipeline through both a command-line interface and a Python API. Two
clustering approaches are supported---Local Estimates Clustering (LEC)
and Extremal Dependence Clustering (EDC)---and every numerical step
is cross-checked against the original R code to make sure results
are consistent across implementations. As a concrete demonstration,
we revisit a published precipitation clustering study for the
AWI-ESM-1-1-LR model ($192\times96$ grid, 1850--2005), running the
full pipeline in Python on a high-performance cluster and obtaining
a spatially consistent cluster structure that is broadly comparable
to the published result.

\bigskip
The manuscript follows the SoftwareX author guidelines.
It contains the mandatory metadata table, descriptions of the three
implemented pipeline flows with worked CLI examples, two result
figures, and a code appendix with self-contained snippets for each
flow. The source code is publicly available on GitHub and permanently
archived on Zenodo.

\bigskip
W.~Kchikach built the software, ran all experiments, and wrote the
paper. T.~Dickhaus and G.~Lohmann guided the scientific direction and
reviewed the manuscript. All authors have read the final version and
agree to this submission. We have no competing interests to declare.

\bigskip
We hope \texttt{weatherisk} will be useful to researchers working at
the intersection of spatial statistics and climate science, and we
look forward to the reviewers' comments.

\bigskip\bigskip
\noindent
Yours sincerely,

\bigskip\bigskip
\noindent
Widad Kchikach\\
(on behalf of all co-authors)

\end{document}
